% 自动生成，勿手改（gen_notes.py）
\chapter*{译校批注}
\phantomsection\addcontentsline{toc}{chapter}{译校批注}
\markboth{译校批注}{译校批注}

本附录汇集正文中以源码注释形式给出的译校批注，按原书顺序排列。每条注明所在原书页码，便于对照原书。批注只写三类内容：易错预警、省略论证的点拨、背景直觉。\par\medskip

\section*{引言}
\addcontentsline{toc}{section}{引言}

\noindent \textbf{（原书 p.v）}\ 日译本此处两处都只写「作用素」，原书英文为 power operations，指的是附录 B 的幂算子（$\psi^{k}$ 一类），与第 3 章讨论的一般算子不同；日译本在此省略了「冪」字，译文按原书英文补足。\par\smallskip

\section*{向量丛}
\addcontentsline{toc}{section}{向量丛}

\noindent \textbf{（原书 p.2）}\ 原书此处写“每个 $e\in E$”，按同构 $E|_Y\cong f^*(E)$ 的含义应读作 $e\in p^{-1}(Y)$；否则 $(p(e),e)$ 不在 $Y\times E$ 中。 \textit{（基本定义）}\par\smallskip

\noindent \textbf{（原书 p.3）}\ 这里“$\Gamma(E)$ 是向量空间”只用到逐点线性结构，与截面空间上可赋的拓扑无关；原书此后也只用其线性结构，不涉及 $\Gamma(E)$ 的拓扑。 \textit{（基本定义）}\par\smallskip

\noindent \textbf{（原书 p.3）}\ 此处“$\Iso(V,W)$ 显然为开集”用的是该集合由 $\det\neq0$ 刻画这一事实；$\dim V\neq\dim W$ 时 $\Iso(V,W)$ 为空集，论断仍成立。 \textit{（基本定义）}\par\smallskip

\noindent \textbf{（原书 p.4）}\ 有限维向量空间之间的 $\Hom$ 上只有一种自然的向量空间拓扑，故这里说的连续性与拓扑的选取无关；到了无限维情形，$T$ 的连续性依赖所取的算子拓扑，这正是上一页所提醒的事。 \textit{（向量丛上的运算）}\par\smallskip

\noindent \textbf{（原书 p.5）}\ 原书此处用 $V$ 表示 $T(E)$ 的任意子集，与上文的向量空间 $V$ 同名，只是局部的临时记法。 \textit{（向量丛上的运算）}\par\smallskip

\noindent \textbf{（原书 p.5）}\ 外幂的记号 $\lambda^i$ 在第 3 章将扩充为一整套 $\lambda$-运算 $\lambda_t(E)=\sum_i\lambda^i(E)t^i$，届时 $\lambda$ 同时充当“运算”与“外幂”两种读法，初读宜先把它固定为外幂。 \textit{（向量丛上的运算）}\par\smallskip

\noindent \textbf{（原书 p.6）}\ 这里“存在开邻域 $U$ 使 $\theta|_U$ 为同构”用的是 p.4 末尾的结论：对平凡丛之间的同态，使某个纤维上为同构的点集是开集。原书只用“そうすると”一句带过。 \textit{（子丛与商丛）}\par\smallskip

\noindent \textbf{（原书 p.7）}\ “配以商拓扑”这一句只交代了集合上的拓扑；$E/F$ 之所以是丛，靠的是上一段刚证明的“子丛局部是直和项”，由此每个 $E|_U/F|_U\cong U\times(V/W)$ 平凡。原书紧接着一句话补上了这一点。 \textit{（子丛与商丛）}\par\smallskip

\noindent \textbf{（原书 p.8）}\ 原书这里的交换图只画了方块，未标出两条水平射各自是什么；上水平射是 $\ker\varphi_x\subset F_x$ 的包含，下水平射是把 $f\in F_x^{**}$ 限制到 $\Coker\varphi_x^*$ 上所得的线性泛函。 \textit{（子丛与商丛）}\par\smallskip

\noindent \textbf{（原书 p.8）}\ 这里把幂等同态称作“投影算子”，是纯代数的叫法；它未必是到子丛的正交投影，因为此时还没有度量。 \textit{（子丛与商丛）}\par\smallskip

\noindent \textbf{（原书 p.8）}\ “两个上半连续函数之和为常值函数，故两者都局部常值”这一步是原书一带而过的关键：上半连续加下半连续才得连续，此处靠的是 $\operatorname{rank}P_x+\operatorname{rank}(1-P_x)\equiv\dim E_x$ 常数这一恒等式。 \textit{（子丛与商丛）}\par\smallskip

\noindent \textbf{（原书 p.9）}\ 页图此处作 $v\in F_x$，但 $P_x$ 的定义域是 $E_x$，右端的 $h_x(v,f_i(x))$ 也要求 $v\in E_x$；应读作 $v\in E_x$，否则该式只说明 $P_x$ 在 $F_x$ 上是恒等映射。 \textit{（子丛与商丛）}\par\smallskip

\noindent \textbf{（原书 p.10）}\ 单位分解的存在性本节直接用，其证明要点是：紧 Hausdorff 空间正规，用 Tietze 扩张定理先造出“贴近 $U_i$ 的”函数，再作归一化求和。原书把它作为拓扑空间论的已知结果引用。 \textit{（紧空间上的向量丛）}\par\smallskip

\noindent \textbf{（原书 p.10）}\ 此处 $t_\alpha$ 只在 $U_\alpha$ 上有定义，故 $S_\alpha$ 必须靠 $\operatorname{supp}p_\alpha\subset U_\alpha$ 才能整体定义；这是单位分解在本节中的第一种用法。 \textit{（紧空间上的向量丛）}\par\smallskip

\noindent \textbf{（原书 p.11）}\ 这一步是本节的技术核心：之所以能从“每一点 $t$ 处同构”跳到“整个 $I$ 上同构类恒定”，靠的是 $I$ 的连通性与同构类函数的局部常值性；后者又只用到了 $Y$ 紧。 \textit{（紧空间上的向量丛）}\par\smallskip

\noindent \textbf{（原书 p.13）}\ 原书 p.13 此处作“$V$ は可縮である”，但论证需要的是 $Y$ 可缩：两个平凡化只相差映射 $Y\to\mathrm{GL}(V)$，要让它与常值映射同伦，靠的正是引理 1.4.7 假设中的 $Y$ 可缩。疑为笔误，正确应为“$Y$ 可缩”。 \textit{（紧空间上的向量丛）}\par\smallskip

\noindent \textbf{（原书 p.15）}\ 这里 $f_t^*(\pi^*E_1\cup_\Phi\pi^*E_2)$ 与 $E_1\cup_{\varphi_t}E_2$ 的同构是逐纤维给出的，原书只用一句“が成り立つ”带过；要点是把 $(y,t)$ 上的纤维分别对应到 $y$ 上两个丛的纤维。 \textit{（紧空间上的向量丛）}\par\smallskip

\noindent \textbf{（原书 p.17）}\ 此处“$(E')^{\perp}\cong E''$”一步省略了细节：完全性给出 $\varphi''$ 在 $(E')^{\perp}$ 上的限制是单射且像为 $E''$，数维数即得同构；这一步依赖每个纤维上的维数相等，而完全序列正是逐纤维的。 \textit{（紧空间上的向量丛）}\par\smallskip

\noindent \textbf{（原书 p.18）}\ 这里的 $G_n(V)$ 是“余维 $n$”的子空间组成的集合，不是“$n$ 维”子空间；后面的 $n=\dim E$ 正是这个意思，因为满射 $\varphi_x\colon V\to E_x$ 的核余维为 $n$。初读极易把它当成通常的 $n$ 维 Grassmann 流形。 \textit{（紧空间上的向量丛）}\par\smallskip

\noindent \textbf{（原书 p.19）}\ 原书 $G_n(V)$ 是 $V$ 的余维 $n$ 子空间全体，不是 $n$ 维子空间全体（定义见原书 p.18）。因此当 $\dim E=n$ 时，$G_n(\C^m)$ 的参数是 $m-n$ 维子空间，也可看成 $n$ 维商空间。若按「$n$ 维子空间」理解定理 1.4.15，维数会对不上。 \textit{（紧空间上的向量丛）}\par\smallskip

\noindent \textbf{（原书 p.20）}\ 这句话只断言 $\Gamma$ 在平凡丛与自由模的子范畴之间是等价；要把等价扩张到全部向量丛，还得靠下一段取投影算子的像所成的子范畴，并配合系 1.4.14。 \textit{（紧空间上的向量丛）}\par\smallskip

\noindent \textbf{（原书 p.21）}\ 这里的 $1$ 是平凡线丛，不是数字；$\mathrm{HOM}(V\otimes V,1)$ 是 $V\otimes V$ 的对偶丛的截面空间，也就是 $V$ 上双线性形式组成的向量空间。前面各节只用 $\Hom$ 记向量空间之间的同态，此处大写的 $\mathrm{HOM}$ 专指丛同态的截面集合。 \textit{（附加结构）}\par\smallskip

\noindent \textbf{（原书 p.22）}\ 自共轭丛的 $T$ 是反线性同构，故 $T^2$ 才是线性的自同态；$T^2=\mathrm{id}_V$ 给出实结构，$T^2=-\mathrm{id}_V$ 给出四元数结构，这两个符号约定在第 2 章讨论 $K$ 理论中的实结构时会反复用到。 \textit{（附加结构）}\par\smallskip

\noindent \textbf{（原书 p.23）}\ 原书 p.23 此处作 $S(X)$，但本节及前面各节都没有定义过 $S(X)$；按上下文，这句的意思应是 $m$ 丛应与 $K$ 群模 $m$ 的对象对应，$S(X)$ 疑为 $K(X)$ 或同类记号的笔误。 \textit{（附加结构）}\par\smallskip

\noindent \textbf{（原书 p.25）}\ 粗体的 $\mathbf{V}_i=X\times V_i$ 是平凡 $G$ 向量丛，斜体的 $V_i$ 是不可约 $G$ 模；上面这个同构把丛的分解化归为表示的分解，粗体与斜体混用时须分清。 \textit{（$G$ 空间上的 $G$-向量丛）}\par\smallskip

\noindent \textbf{（原书 p.26）}\ 平均化算子要除以 $|G|$，所以在特征不整除 $|G|$ 的域上才可逆；本书只在复系数上工作，有限群这一限制正是为了让平均化恒可行。若换成紧群，则应把求和换成 Haar 积分。 \textit{（$G$ 空间上的 $G$-向量丛）}\par\smallskip

\noindent \textbf{（原书 p.28）}\ 这个 $B$ 就是交叉积 $A\rtimes G$（扭群环）：乘法里 $g$ 先作用到系数 $a'$ 上再相乘，因此 $B$ 既不是 $A$ 与群环 $\C[G]$ 的张量积，也不是直积；称它为「扭」群环正是指这种系数被 $G$ 扭一下的乘法。 \textit{（$G$ 空间上的 $G$-向量丛）}\par\smallskip

\section*{K 理论}
\addcontentsline{toc}{section}{K 理论}

\noindent \textbf{（原书 p.30）}\ 普遍性只保证 $\alpha\colon A\to K(A)$ 是半群同态，并不保证它是单射；$A$ 若没有消去律，不同的元素可能在 $K(A)$ 中有同一个像。对 $A=\Vect(X)$ 也如此，这正是下文必须引入稳定等价的原因。 \textit{（定义）}\par\smallskip

\noindent \textbf{（原书 p.30）}\ 带下划线的 $\underline{n}$ 专指 $n$ 维平凡向量丛，与整数 $n$ 区分；后文还会出现 $[\underline{n}]$ 与 $\underline{n}$ 混用之处，读者须按上下文判断哪个是 $K(X)$ 中的类。 \textit{（定义）}\par\smallskip

\noindent \textbf{（原书 p.32）}\ 这里的 $P^0$ 是单位圆盘丛（长度不超过 1 的部分），$P^\infty$ 是补上无穷远截面后的外部，两者都是 $P(L\oplus\underline{1})$ 的闭子集；$S$ 是它们的公共边界。零截面落在 $P^0$ 内，无穷远截面是 $P^\infty$ 的「尖端」。 \textit{（周期性定理）}\par\smallskip

\noindent \textbf{（原书 p.33）}\ 原书此句前半作 $a_k\in L^{-k}$，紧接的后一句却写作 $a_k\in\Gamma(L^{-k})$；按上下文，$a_k$ 是 $L^{-k}$ 的截面，前一处漏排了 $\Gamma$。译文照录原文，读者按截面理解即可。 \textit{（周期性定理）}\par\smallskip

\noindent \textbf{（原书 p.34）}\ 原书把 Hom 空间的截面集记作 $\Gamma(\pi^{*}(E^0),\pi^{*}(E^\infty))$，但 $\Gamma$ 只取一个变元，这个记号与 p.20 的 $\mathrm{HOM}$ 用法不一致；按紧接的 $f\in\mathrm{ISO}(\cdots)$ 判断，此处指的是丛同态 $\Hom(\pi^{*}(E^0),\pi^{*}(E^\infty))$ 的截面。 \textit{（周期性定理）}\par\smallskip

\noindent \textbf{（原书 p.34）}\ 原书 p.34 把 $f_\infty$ 的靶印作 $\pi_0^*(L)$，但 $\pi_0$ 是 $P^0$ 上的投影，$f_\infty$ 定义在 $P^\infty$ 上，此处应为 $\pi_\infty^*(L)$，疑为笔误。译文照录原文。 \textit{（周期性定理）}\par\smallskip

\noindent \textbf{（原书 p.35）}\ 这里 $f$ 是紧空间 $S$ 上的连续截面，Fourier 系数用逐纤维的围道积分定义；系数 $a_k$ 落在 $\Gamma(\Hom(L^k\otimes E^0,E^\infty))$ 中，是因为 $f_xz_x^{-k-1}dz_x$ 沿 $S_x$ 积分后在 $L_x^{-k}$ 中取值，再与 $E^0_x$ 张量并作用到 $E^\infty_x$。 \textit{（周期性定理）}\par\smallskip

\noindent \textbf{（原书 p.38）}\ 原书 p.38 此处作「命題 2.2.7」，但这一段正在证明命题 2.2.7，自引不合逻辑；由所引内容（$\mathcal{L}^n(p)$ 与 $p\oplus1$ 界定同构的丛）判断，应为命题 2.2.5。译文照录原文。 \textit{（周期性定理）}\par\smallskip

\noindent \textbf{（原书 p.39）}\ 易错预警：原书 p.39 定义的两个积分顺序相反，$Q^0$ 用 $p_x^{-1}dp_x$，$Q^\infty$ 用 $dp_xp_x^{-1}$。对非交换的自同态族这两者并不相等，正是积分顺序的差异使 $pQ^0=Q^\infty p$ 成立，抄录时不可把两个积分写成同一个。 \textit{（周期性定理）}\par\smallskip

\noindent \textbf{（原书 p.42）}\ 易错预警：上两式里的 $[H^{-1}]$ 就是 $[H]^{-1}$，$H^{-1}$ 是 $H$ 的对偶线丛，$[H][H^{-1}]=[1]$；因此 $[H^{-1}]-[1]$ 才能与前面的 $([H]^{n-1}-[H]^n)$ 对应。原书在 $[H]$ 与 $[H^{-1}]$ 之间来回换写，读时按同一个元素处理。 \textit{（周期性定理）}\par\smallskip

\noindent \textbf{（原书 p.43）}\ 易错预警：$\nu_n(f)$ 的系数取 $V_{2n}$ 而不是 $V_n$，下标 $2n$ 来自 $p_n=z^nf_n$ 的次数上界 $2n$；后文出现的 $V_{2n+1}$、$V_{2n+2}$ 也据此递推。 \textit{（周期性定理）}\par\smallskip

\noindent \textbf{（原书 p.45）}\ 省略论证的点拨：定理 2.3.1 没有单列证明，它就在上一段里。只要把 §2.2 的每一步（$P(L\oplus1)$ 的分解、$z$ 与 $H$ 的截面、衔接函数的同伦）逐字换成 $G$ 等变版本即可，所有构造都带 $G$ 作用。 \textit{（$K_G(X)$）}\par\smallskip

\noindent \textbf{（原书 p.45）}\ 记号提醒：原书 p.45 把模写作 $K_G(X)[t]/((t[L]-1)(t-1))$，而 p.42、p.43 写作 $K(X)[t]/((t-1)([L]t-1))$。因 $t$ 与 $[L]$ 可交换，两式生成同一理想，并非笔误，比较时不要误以为差一个符号。 \textit{（$K_G(X)$）}\par\smallskip

\noindent \textbf{（原书 p.46）}\ 易错预警：$K(X,Y)$ 的定义是 $\tilde K(X/Y)$，而不是 $K(X)$ 对 $K(Y)$ 的商。只有 $Y\to X$ 是好对（例如 $Y$ 是 $X$ 的形变收缩）时，二者才一致；一般情形下用 $K(X/Y)$ 才是对的。 \textit{（K 群的一般上同调论性质）}\par\smallskip

\noindent \textbf{（原书 p.46）}\ 次数约定：定义 2.4.1 只给出负次数 $K^{-n}$（$n\ge0$），正次数要靠 Bott 周期性 $\tilde K(S^2X)\cong\tilde K(X)$ 反向延拓，即令 $K^{n}=K^{n-2}$。这是后文 $K^*(X)$ 记号的基础，不要以为漏了正次数。 \textit{（K 群的一般上同调论性质）}\par\smallskip

\noindent \textbf{（原书 p.47）}\ 点拨：$CX/X$ 与 $SX$ 只差把子空间 $I$ 收缩成一点，故用引理 1.4.7（紧空间上的同伦等价给出 $K$ 群的同构）即得 $K(SX)\cong K(CX/X)$。这一步是把相对 $K$ 群与约化悬垂接起来的关键，后面的长正合列全靠它。 \textit{（K 群的一般上同调论性质）}\par\smallskip

\noindent \textbf{（原书 p.48）}\ 点拨：证明中 $i^*j^*=0$ 的理由是，$j^*$ 之后复合的 $ji$ 通过 $(Y,Y)$ 分解，而 $K(Y,Y)=\tilde K(Y/Y)=0$。这是“平凡群”三个字的实际含义。 \textit{（K 群的一般上同调论性质）}\par\smallskip

\noindent \textbf{（原书 p.52）}\ 点拨：$\delta$ 定义为 $m^*\theta^{-1}$，先由 $\theta$ 把 $K(X\cup CY,X)$ 送到 $\tilde K^{-1}(Y)$，再用 $m^*$ 拉回。两个映射的顺序不能颠倒，$\delta$ 的符号完全由这个复合决定。 \textit{（K 群的一般上同调论性质）}\par\smallskip

\noindent \textbf{（原书 p.52）}\ 点拨：图式 (D) 的左列是锥 $C_2$、右列是锥 $C_1$，交换这两个锥的因子恰好差一个符号；这正是引理 2.4.5 中 $(T\wedge1)^*y=-y$ 的来源，也是作者反复强调“除符号之外”的原因。 \textit{（K 群的一般上同调论性质）}\par\smallskip

\noindent \textbf{（原书 p.52）}\ 原书 p.52 此处作 $y\in\tilde K(SX)$，与同一陈述中的 $T\wedge1\colon SY\to SY$ 及下文的用法不符，疑为 $\tilde K(SY)$ 之误。正文照原书抄录，请读者按 $SY$ 理解。 \textit{（K 群的一般上同调论性质）}\par\smallskip

\noindent \textbf{（原书 p.55）}\ 点拨：6 项正合列画成一个循环，两个竖直箭头都是 $\delta$：右边是 $K^0(Y)\xrightarrow{\delta}K^1(X,Y)$，左边是 $K^1(Y)\xrightarrow{\delta}K^0(X,Y)$。把 $K^n$ 与 $K^{n-2}$ 同一化之后，$\delta$ 才能首尾相接成环。 \textit{（K 群的一般上同调论性质）}\par\smallskip

\noindent \textbf{（原书 p.56）}\ 点拨：定理 2.5.1 的证明只用到一件事——$S^{n-1}$ 不是 $D^n$ 的收缩核。而这一点由 $\tilde K^*(D^n)=0$、$\tilde K^*(S^{n-1})\ne0$ 得出：收缩核会给出 $K^*$ 上的分裂单射，与两者不同构矛盾。 \textit{（$K^*(X)$ 的计算）}\par\smallskip

\noindent \textbf{（原书 p.57）}\ 次数约定：命题 2.5.2 中的 $K^1(X)=0$ 依赖把 $K^n$ 与 $K^{n-2}$ 同一化的约定；这里的 $K^1$ 就是奇数次群 $K^{-1}$，不是另一个独立的群。 \textit{（$K^*(X)$ 的计算）}\par\smallskip

\noindent \textbf{（原书 p.59）}\ 原书用 $\hat\otimes$ 表示外部张量积、用 $\otimes$ 表示普通张量积：$E\hat\otimes F$ 定义在积空间 $X\times Y$ 上，而 $E\otimes F$ 只在同一底空间上才有意义，它是 $\Delta^*$ 的拉回。初读时容易把两者混为一谈。 \textit{（$K^*(X,Y)$ 中的乘法）}\par\smallskip

\noindent \textbf{（原书 p.60）}\ 这里的 $\beta$ 就是 Bott 元素。乘 $\beta$ 给出的平移同构 $K^{-n}\cong K^{-n-2}$ 正是 Bott 周期性的实质内容；而 $K^*=K^{\#}/(1-\beta)$ 把这条周期性直接做成环的商，故 $K^*$ 只在模 2 意义下分次。 \textit{（$K^*(X,Y)$ 中的乘法）}\par\smallskip

\noindent \textbf{（原书 p.63）}\ 定义 2.6.3 说的是“自然变换”，即 $\chi_n$ 对所有空间对 $(X,A)$ 同时给出，并与拉回相容；只在某个固定对上有交替和并不足以确定它。$\chi_n$ 的存在性要到引理 2.6.8 才证明。 \textit{（$K^*(X,Y)$ 中的乘法）}\par\smallskip

\noindent \textbf{（原书 p.66）}\ 引理 2.6.11 的技巧是全节的关键：把复形 $\cdots\to E_{2i}\to E_{2i-1}\to\cdots$ 的偶、奇两组直和分别记作 $F_0,F_1$，再把 $\alpha$ 与其 Hermite 共轭拼成一个映射 $\beta\colon F_1\to F_0$。由正合性，$\beta$ 是 $A$ 上的同构，于是任意长的复形都压缩为 $\mathcal{L}_1$ 中的一个元素。 \textit{（$K^*(X,Y)$ 中的乘法）}\par\smallskip

\noindent \textbf{（原书 p.69）}\ 这个复形的映射就是外积代数上乘 $v$，即 Koszul 复形。$v\ne0$ 时乘 $v$ 在 $\Lambda^*(V)$ 上可逆（收缩映射给出其逆），故复形正合；作者省略的正是这一句。它也是全书后面 Thom 类构造的原型。 \textit{（$K^*(X,Y)$ 中的乘法）}\par\smallskip

\noindent \textbf{（原书 p.70）}\ 原书在空间对的 $K(B(V),S(V))$ 与商空间的约化群 $\tilde K(X^V)$ 之间不加区分地换用：二者经 $B(V)/S(V)\simeq X^V$ 等同。当 $V$ 的纤维是 $\C^n$ 时 $\tilde K(X^V)$ 只由维数 $2n$ 决定，故性质 (A) 实际上是在说“逐纤维看 Thom 类就是标准的 Bott 生成元”。 \textit{（$K^*(X,Y)$ 中的乘法）}\par\smallskip

\noindent \textbf{（原书 p.73）}\ 点拨：命题 2.7.1 的关系式 $\prod_i([L_i][H]-1)=0$ 与命题 2.5.3 的 $\prod_i(t-[L_i^*])$ 是同一个关系：令 $t=[H]^{-1}=[H^*]$ 再乘以 $[H]^n$ 即可互推。抄录时注意两式中 $1$ 与 $[H]$ 的位置不同。 \textit{（Thom 同构）}\par\smallskip

\noindent \textbf{（原书 p.75）}\ 易错预警：$R(\Z/2)=\Z[\rho]/(\rho^2-1)$ 中 $\rho$ 是自然表示的特征标，$\sigma=\rho-1$ 对应一维非平凡表示；关系 $\sigma^2=-2\sigma$ 表明 $\Z[\rho]/(\rho^2-1)$ 的两个生成元之和与差都给出 $\Z$ 的直和项。$\lambda_{-1}[E]=(-\sigma)^n$ 的形式决定下文 $2^{n-1}$ 的阶，符号不能丢。 \textit{（Thom 同构）}\par\smallskip

\noindent \textbf{（原书 p.77）}\ 点拨：定理 2.7.8 的论证骨架是“局部同构 $\Rightarrow$ 两开集之并仍良好 $\Rightarrow$ 由紧性得整体”。中间那一步用对 $(V,V_2)$ 的正合列（Mayer–Vietoris）把 $K^*(V)$ 夹在中间；$M^*$ 无挠保证 $\otimes M^*$ 保持单射，此条件不可省。 \textit{（Thom 同构）}\par\smallskip

\noindent \textbf{（原书 p.78）}\ 点拨：分裂原理之所以够用，是因为 $\pi^*$ 是单射：一个等式在 $K^*(X)$ 中成立当且仅当它在 $K^*(F)$ 中成立，而 $\pi^*(E_i)$ 已分解为线丛之和，命题 2.7.1、2.7.2 便可逐项使用。这也解释了作者为何先算可分解情形。 \textit{（Thom 同构）}\par\smallskip

\noindent \textbf{（原书 p.79）}\ 原书 p.79 把 Künneth 公式中的 $K^j(Y)$ 印成了 $K^j(X)$（两处都是）。按公式的含义应为 $K^j(Y)$：左端是 $K^*(X)\otimes K^*(Y)\to K^*(X\times Y)$ 的 Künneth 分解。正文照原书抄录，请按 $K^j(Y)$ 理解。 \textit{（Thom 同构）}\par\smallskip

\noindent \textbf{（原书 p.80）}\ 次数约定：原书此处写 $K^{-1}(X)=K(SX)$；按定义 2.4.1，$K^{-1}(X)=\tilde K(S(X^+))$，故右端应理解为约化群 $\tilde K(SX)$。$K(SX)$ 还多一个 $\Z$ 的直和项，读时按约化群处理。 \textit{（Thom 同构）}\par\smallskip

\noindent \textbf{（原书 p.81）}\ 点拨：$\lambda^i$ 到 $U(n-1)$ 的限制为 $\mu^i\otimes\mu^{i-1}$，这来自外积代数的分裂 $\Lambda^*(\C^n)=\Lambda^*(\C^{n-1})\oplus e_n\Lambda^*(\C^{n-1})$：含 $e_n$ 的那一半把次数升 1，正是 $\mu^{i-1}$ 乘以 $\mu^1$ 的效果。这一步是归纳法能启动的关键。 \textit{（Thom 同构）}\par\smallskip

\section*{算子}
\addcontentsline{toc}{section}{算子}

\noindent \textbf{（原书 p.83）}\ 这里要求的是自然变换，而不是对各个空间任意指定的自同态：$F_X f^*=f^*F_Y$ 必须对一切连续映射 $f$ 成立。只在一个空间上给定的同态一般拼不成算子；本节构造的 $\lambda_t$、$\gamma_t$ 都自动满足这一自然性。 \textit{（外幂算子）}\par\smallskip

\noindent \textbf{（原书 p.85）}\ 此处的「正」只是说 $x$ 落在 $\Vect(X)$ 的像中，并非来自 $K(X)$ 的任何序结构；$K(X)$ 上没有与 $\oplus$ 相容的全序，切勿把 $x\ge 0$ 读成实数大小。 \textit{（外幂算子）}\par\smallskip

\noindent \textbf{（原书 p.86）}\ 原日译本 p.85 此处印作 $\rank E_x$，同一段下文与全节各处则作 $\rank E$。这里的 $\rank E$ 是一个丛（各连通分支上纤维维数为 $\dim E_x$ 的平凡丛），不是秩函数，下标 $x$ 疑为衍字；后文 $E\mapsto\rank E$ 才是那个环自同态。 \textit{（外幂算子）}\par\smallskip

\noindent \textbf{（原书 p.87）}\ $\gamma_t$ 只是 $\lambda_t$ 的变量替换，并不产生新的信息；改用 $\gamma^i$ 的理由在命题 3.1.1：$\gamma$-类在 $K_1(X)$ 上幂零，因而形式幂级数可作用于任意 $x\in K_1$，这正是定理 3.1.7 的出发点。 \textit{（外幂算子）}\par\smallskip

\noindent \textbf{（原书 p.87）}\ 证明中「$\dim E=n$」省了一步：由 $x+n=[E]$ 与 $\rank x=0$ 得 $\rank E=n$，即 $E$ 在各连通分支上的纤维维数是 $n$，故 $\dim E=n$；否则后面 $\lambda^i(E)=0$（$i>n$）不能直接用。 \textit{（外幂算子）}\par\smallskip

\noindent \textbf{（原书 p.88）}\ 界 $n(n-1)$ 来自「$n$ 个变量、每个指数不超过 $n-1$」时可能达到的最高次数；证明中「被某个 $x_i^n$ 整除」用的是抽屉原理：次数一旦超过 $n(n-1)$，必有某个指数不小于 $n$。 \textit{（外幂算子）}\par\smallskip

\noindent \textbf{（原书 p.89）}\ 作者在此明确划出两条路线：引理 3.1.3 是「代数的分裂原理」，只依赖对称函数恒等式；§2.7 的位相分裂原理依赖 Bott 周期性。若用后者去证本节的形式结论，就会绕回第 2 章的深刻定理，形成循环论证。 \textit{（外幂算子）}\par\smallskip

\noindent \textbf{（原书 p.92）}\ $K(Y_{n,m})$ 一行中的 $x_i$ 不是超平面线丛的类本身，而是 $[H_i]-1$；正因如此才有 $\gamma_t(x_1+\cdots+x_n)=\prod(1+x_it)$。若误把 $x_i$ 当作 $[H_i]$，这条公式会差一个因子。 \textit{（外幂算子）}\par\smallskip

\noindent \textbf{（原书 p.94）}\ 这里的 $\prod$ 在原书中表示各连通分支的不交并，而非拓扑积；$H^0$ 在每个连通分支上取 $\Z$，故不交并给出直积 $\Z^n$，两种读法在此恰好一致。 \textit{（外幂算子）}\par\smallskip

\noindent \textbf{（原书 p.96）}\ 这里用的是 $\lambda_{-t}$ 而非 $\lambda_t$，并且前面带 $-t\frac{d}{dt}$：这两处符号配合起来，才使线丛情形给出 $\psi^k(x)=x^k$（命题 3.2.1 第 2 项）；写成 $\lambda_t$ 会把符号弄反。 \textit{（Adams 算子）}\par\smallskip

\noindent \textbf{（原书 p.96）}\ 原日译本 p.96 此处（以及本命题第 2 项、证明第二段）印作大写的 $X$：$X$ 是空间记号，此处应指元素 $x$ 为线丛的类，大写 $X$ 疑为笔误；正文照印。 \textit{（Adams 算子）}\par\smallskip

\noindent \textbf{（原书 p.97）}\ 为何 $x^2$ 必是 $y$ 的整数倍：$\tilde K(X_f)$ 以 $x,y$ 为自由 Abel 基，而 $i^*(x^2)=(i^*x)^2=0$ 说明 $x^2$ 落在 $\tilde K(X_f,S^{2n})$ 中，这个子群由 $y=j^*(u)$ 生成，故 $x^2=Hy$。Hopf 不变量的整性是这一构造的全部要点。 \textit{（Adams 算子）}\par\smallskip

\noindent \textbf{（原书 p.99）}\ 这个乘积公式的来历：命题 3.2.4 说每个因子 $\psi^k-k^n$ 把 $\gamma$ 滤过至少提高 1；从 $F_\gamma\ge0$ 出发乘上 $m=0,1,\dots,n$ 这 $n+1$ 个因子后落在 $K_{n+1}^\gamma(X)=0$ 中，故整体为零算子。 \textit{（Adams 算子）}\par\smallskip

\noindent \textbf{（原书 p.100）}\ 这里把 $\psi^k$ 的特征空间直接命名为 $H^{2i}(X;\Q)$，只给出偶数次的上同调；奇数次靠 $\psi^k$ 看不到，才用悬挂 $H^{2m+1}(X;\Q)=H^{2m+2}(SX^+;\Q)$ 人为定义。故这套「Betti 群」与通常奇异上同调的一致性并非平凡事实。 \textit{（Adams 算子）}\par\smallskip

\noindent \textbf{（原书 p.101）}\ 原书中公式编号与陈述编号共用「节.序号」计数器：命题 3.2.7 之后直接是公式 (3.2.8)，再接命题 3.2.9。抄录时不要误补出一个「命题 3.2.8」。 \textit{（Adams 算子）}\par\smallskip

\noindent \textbf{（原书 p.103）}\ 限制到 $S^{2n}$ 时 $\tilde K(S^{2n})$ 中 $x^{n+1}=0$，故 $i^*(x^i)=0$（$i>n$）；同理 $i^*(y)$ 与 $i^*(x)$ 在 $K(S^{2n})$ 中相等，这正是 $\lambda=1$ 的来源。 \textit{（Adams 算子）}\par\smallskip

\noindent \textbf{（原书 p.104）}\ 定义只用到球面丛 $S(E)$ 与投影 $S(E)\to X$，与向量丛的加法结构无关；因此 $E$ 与 $E\oplus V$ 一般并不等价，必须再做稳定化才得到群。 \textit{（群 $J(X)$）}\par\smallskip

\noindent \textbf{（原书 p.105）}\ 稳定同伦型在悬垂下不变，而 $S^{2n}(X^+)$ 与 $X^+$ 稳定同伦等价，所以这里可以用 $X^+$ 代替平凡丛的 Thom 复形。 \textit{（群 $J(X)$）}\par\smallskip

\noindent \textbf{（原书 p.105）}\ 这里 $\lambda_E\in\tilde K(X^E)$ 是 Thom 类；$\tilde K(X^E)$ 是秩 1 的自由 $K(X)$ 模，所以 $\psi^k(\lambda_E)$ 必是 $\lambda_E$ 的 $K(X)$-倍数，$\rho^k(E)$ 由此唯一确定。 \textit{（群 $J(X)$）}\par\smallskip

\noindent \textbf{（原书 p.108）}\ 由 $[L]^2=1$ 得 $x^2=([L]-1)^2=2-2[L]=-2x$，这是把 $(1+x/3)^m$ 的所有高次项降为一次的根据。 \textit{（群 $J(X)$）}\par\smallskip

\noindent \textbf{（原书 p.108）}\ $x$ 的阶为 $2^n$，故 $m$ 只须按模 $2^n$ 考虑；在 $\Z/2^n$ 中满足 $2^{n-1}\mid m$ 的只剩 $m\equiv0$ 与 $m\equiv2^{n-1}$，这就是“核的阶至多 2”的来源。 \textit{（群 $J(X)$）}\par\smallskip

\section*{Fredholm 算子的空间}
\addcontentsline{toc}{section}{Fredholm 算子的空间}

\noindent \textbf{（原书 p.111）}\ 定义中两个维数都必须有限，指标才有意义；只要求其中一个有限并不够，因为 $\dim\Ker T$ 与 $\dim\Coker T$ 的差在扰动下才稳定。另注意此处的 index 是整数，与第 2 章中 $\operatorname{index}:[X,\cF]\to K(X)$ 的群值指标是同一概念的不同层级的推广。\par\smallskip

\noindent \textbf{（原书 p.112）}\ 这里是半群 $[X,\cF]$ 到群 $K(X)$ 的同构；$\cF$ 本身不连通，$[X,\cF]$ 只在同伦意义下成半群，故定理的“自然同构”中同构一词指半群同构。\par\smallskip

\noindent \textbf{（原书 p.114）}\ (ii) 中的拓扑是 $X\times H$ 的商拓扑，并非各纤维上欧氏拓扑的简单并；这里“向量丛”的全部内容就是这一步局部平凡性，$\cA$ 上的范数拓扑恰好保证了 $\varphi_S$ 对 $S$ 连续。\par\smallskip

\noindent \textbf{（原书 p.118）}\ $\cA^*$ 可缩是无穷维 Hilbert 空间特有的现象，与有限维时酉群 $U(n)$ 的连通分支结构完全不同；正是这一可缩性使定理 A.1 的分类空间可以取 $\cF$ 而不必另作修正，Kuiper 定理是整篇附录的技术支柱。\par\smallskip

\section*{K 理论中的幂算子}
\addcontentsline{toc}{section}{K 理论中的幂算子}

\noindent \textbf{（原书 p.121）}\ 这里的 $K_q(X)$ 是滤过的第 $q$ 层，即限制到 $(q-1)$ 维骨架上为零的丛类；下标 $q$ 越大子群越小，与“$q$ 次”的直觉相反。下文 $GK(X)$ 取商 $K_{2q}(X)/K_{2q-1}(X)$，正是要把只有滤过的 $K(X)$ 强行分次化。\par\smallskip

\noindent \textbf{（原书 p.122）}\ 下文反复出现的“算子”，指的是函子 $K$ 到自身的自然变换：对每个空间 $X$ 给一个映射 $K(X)\to K(X)$，并与拉回相容。这一点在 p.138 才正式写成 $\operatorname{Op}(K)$，此处先借用其意，初读宜先记住“算子 = 自然的 K 群自映射”。\par\smallskip

\noindent \textbf{（原书 p.124）}\ 上面的分解依据 $S_k$ 的半单性（有限群表示完全可约）：$\pi(E)=\Hom_{S_k}(V_\pi,E^{\otimes k})$ 是 $V_\pi$ 在 $E^{\otimes k}$ 中的重数空间。当 $V_\pi$ 取平凡表示时得对称幂 $\sigma^k$，取符号表示时得外幂 $\lambda^k$，这两个极端情形正是下文挑出来单独研究的。 \textit{（张量积与对称群）}\par\smallskip

\noindent \textbf{（原书 p.125）}\ 左端的 $\operatorname{Trace}_{S_k}(T^{\otimes k})$ 并不是线性映射 $T^{\otimes k}$ 的普通迹（那会等于 $\operatorname{Trace}(T)^k$），而是把 $T^{\otimes k}$ 按 $S_k$ 的不可约表示分块后逐块取迹；正因如此它才落在 $\operatorname{Sym}_k\otimes R(S_k)$ 里，并由函数 $g\mapsto\operatorname{Trace}(gT^{\otimes k})$ 唯一决定。 \textit{（张量积与对称群）}\par\smallskip

\noindent \textbf{（原书 p.128）}\ 证明的关键是秩的比对：$\operatorname{Sym}_k[t_1,\ldots,t_n]$ 的秩等于 $k$ 的分割数，$R_*(S_k)$ 的秩等于 $S_k$ 的共轭类数，而 $S_k$ 的共轭类与 $k$ 的分割一一对应。先证满射、再证秩相等，即得同构；这是有限自由交换群证明同构的常用套路。 \textit{（张量积与对称群）}\par\smallskip

\noindent \textbf{（原书 p.132）}\ 原书 p.131 的式子 $\psi^k=Q_k(\lambda^1,\ldots,\lambda^k)=(-1)^{k-1}k\lambda_k+$ 复合项 把 $\lambda$ 写成下标 $\lambda_k$，与本节其余各处的外幂 $\lambda^k$（上标）不一致；按牛顿恒等式，此处应为 $(-1)^{k-1}k\lambda^k$，疑为笔误。 \textit{（张量积与对称群）}\par\smallskip

\noindent \textbf{（原书 p.132）}\ 等式 $\Lambda_{-1}(M)(g)=\det(1-g_M)$ 来自外幂的特征标公式：$\sum_i(-1)^i\operatorname{Trace}\Lambda^i(g_0)=\det(1-g_0)$。因此 $g_M$ 有特征值 1 时右边为零，这正是“$\Lambda_{-1}(M)$ 在复合类上消失”的依据；而在长度 $k$ 的循环置换上，$M$ 的特征值为除 1 以外的全部 $k$ 次单位根，故 $\prod_{i=1}^{k-1}(1-\zeta^i)=k$。 \textit{（张量积与对称群）}\par\smallskip

\noindent \textbf{（原书 p.135）}\ 这里要求 $X$ 是平凡 $G$ 空间，于是 $g$ 覆盖的底空间映射是恒等，条件 (ii) 自动成为“每个 $g$ 给出 $E$ 的一条纤维自同构”，$G\to\Aut E$ 的表述才有意义。对非平凡 $G$ 空间，$G$ 向量丛只是等变丛，并没有这种同态。 \textit{（K 理论中的算子）}\par\smallskip

\noindent \textbf{（原书 p.137）}\ 交换 $p$ 次与 $q$ 次因子时乘 $(-1)^{pq}$ 是 Koszul 符号规则，正是为了让置换作用成为复形的自同构：否则 $d^{\otimes k}$ 与重排不对易，下面证明中“$Q$ 作为复形非轮状”就不成立。这一步是易错点。 \textit{（K 理论中的算子）}\par\smallskip

\noindent \textbf{（原书 p.140）}\ $Q_k$ 是 $k$ 次 Newton 多项式，它把初等对称函数 $e_1,\dots,e_k$ 与幂和 $p_k=\sum t_i^k$ 互化；Adams 正是用它由外幂 $\lambda^k$ 反解出 $\psi^k$。所以“$\psi^k$ 由 $\lambda^1,\dots,\lambda^k$ 表示”不是定义选择，而是 Newton 恒等式的直接推论。 \textit{（K 理论中的算子）}\par\smallskip

\noindent \textbf{（原书 p.143）}\ 这一套与 §2.6 的 $\mathcal{C}_n$、$\mathcal{L}_n$ 是同一台机器的等变版本，区别只在：这里允许任意长的复形（不固定长度 $n$），且所有映射须与 $G$ 的作用可交换。基本对象 $P$ 的加法就是 §2.6 中的“基本列”。 \textit{（外部张量幂）}\par\smallskip

\noindent \textbf{（原书 p.145）}\ “$q$ 切片”是日译本对骨架（skeleton）的译名，本书一律称“$q$ 骨架”。注意 $K_q(X)$ 由限制映射 $K(X)\to K(X_{q-1})$ 的核定义：$q$ 越大条件越强、$K_q$ 越小，故这是“余维 $\ge q$”型的滤过，与按骨架递增的拓扑滤过方向相反。 \textit{（算子与滤过）}\par\smallskip

\noindent \textbf{（原书 p.146）}\ $K(B_G)$、$K(X_G)$ 的定义是沿骨架取正向极限，故 (B.4.1) 中出现的 $K(X\times B_{S_k})$ 实际应读作 $\varinjlim_q K(X\times B_{S_k,q})$；书中此后不再重复这层意思。 \textit{（算子与滤过）}\par\smallskip

\noindent \textbf{（原书 p.148）}\ 证明中 $X^{k-1}X_{q-1}$ 是简写，指乘积空间 $X^{k-1}\times X_{q-1}$；$A$ 即 $B_{S_k}$ 的万有覆盖。 \textit{（算子与滤过）}\par\smallskip

\noindent \textbf{（原书 p.149）}\ 原书 p.149 此处印作 $\lambda^i(X)\lambda^j(x)$ 与 $\lambda^k(X)$，但同一页的证明随即写作 $\lambda^k(x)$，且 $\lambda$ 运算作用在 $K$ 群元素上，故疑为笔误，应读作 $\lambda^i(x)\lambda^j(x)$ 与 $\lambda^k(x)$。 \textit{（算子与滤过）}\par\smallskip

\noindent \textbf{（原书 p.150）}\ $\widehat{R(\Z/k)}$ 表示表示环 $R(\Z/k)$ 对 $I(\Z/k)$-进拓扑的完备化；$I(\Z/k)$ 由增广理想（即 $R$ 中维数部分为零的元素）生成，$I(\Z/k)^q$ 是它的 $q$ 次幂。 \textit{（算子与滤过）}\par\smallskip

\noindent \textbf{（原书 p.152）}\ 推论 B.5.2 的 $k$ 由 $k+1\ge n/(p-1)$ 的最小整解决定，故 $k=\lceil n/(p-1)\rceil-1$；它与定理 B.5.3 证明中出现的 $k$ 是同一个量，注意两处下标不要混。 \textit{（素数阶循环群的情形）}\par\smallskip

\noindent \textbf{（原书 p.153）}\ 原书 p.153 此处引“系 B.5.2”，但 $\theta^p(x)$ 的可除性估计实际来自定理 B.5.3；系 B.5.2 只给出像的阶为 $p^k$。照原书抄录，此处的引用宜读作定理 B.5.3。 \textit{（素数阶循环群的情形）}\par\smallskip

\noindent \textbf{（原书 p.154）}\ $\theta^p(x)$ 是 $P^p(x)=\theta^p(x)\otimes(N-p)$ 中的因子，$\psi^p$ 是 Adams 算子，二者由 $\psi^p(x)=x^p-p\theta^p(x)$ 相联系；命题 B.5.6 把 $\psi^p(x)$ 按 $p$ 的幂次展开，正是后面比较 Steenrod 作用素的入口。 \textit{（素数阶循环群的情形）}\par\smallskip

\noindent \textbf{（原书 p.155）}\ $H^*(X;\Z)\Rightarrow K^*(X)$ 中的双箭头指 Atiyah–Hirzebruch 谱序列；文中“由 $a$ 代表”是说 $a$ 是 $E_2$ 页上对应 $x$ 的类，而不是 $x$ 的某个具体代表元，故这类断言只能逐次读取。 \textit{（素数阶循环群的情形）}\par\smallskip

\noindent \textbf{（原书 p.156）}\ (B.6.3) 中 $(c\otimes\cdots\otimes c)$ 是 $k$ 个 $c$ 的张量积，$\otimes_\Gamma$ 再用群环 $\Gamma$ 与 $1$ 作用，等于把 $S_k$ 的对称化一次并入；这也是为何 $\nu$ 只在 $X_{S_k}$ 一侧出现。 \textit{（与上同调算子的关系）}\par\smallskip

\noindent \textbf{（原书 p.157）}\ $p=2$ 时约定 $P^i_2=Sq^{2i}$，只取偶数次 Steenrod 平方；这与奇素数时 $P^i_p$ 的次数 $2i(p-1)$ 相合，若误写成 $Sq^i$ 则次数全错。 \textit{（与上同调算子的关系）}\par\smallskip

\noindent \textbf{（原书 p.157）}\ 此处的 $\eta\in H^2(\Z/p;\Z/p)$ 是上同调的标准生成元，与 B.5 中 $\eta=\rho-1\in R(\Z/p)$ 只是同名的不同对象，两者不可混用；书上不给新记号，读时须按所在环判断。 \textit{（与上同调算子的关系）}\par\smallskip

\noindent \textbf{（原书 p.160）}\ “$p$ 整元”是把 $K$ 群中的整除性搬到 $H^*(X;\Q)$ 的唯一通道：$H^*(X;\Q)$ 里除法不受限制，只有把分母选成与 $p$ 互素，模 $p$ 约化才有意义。 \textit{（与 Chern 特征标的关系）}\par\smallskip

\section*{K 理论与实结构}
\addcontentsline{toc}{section}{K 理论与实结构}

\noindent \textbf{（原书 p.164）}\ 原书 p.164 此式左端印作大写的 $\tau(X)$，右端却是 $\bar x$；按下文“把对合像复共轭那样写”的用法，应为 $\tau(x)=\bar x$，疑为排版笔误。 \textit{（实空间的范畴）}\par\smallskip

\noindent \textbf{（原书 p.165）}\ 这段话是理解全篇的关键：同一个复向量丛 $E$ 配上“线性”的对合就落在 $\Z/2$ 空间范畴里，配上“反线性”的对合才落在实空间范畴里；$KR$ 与 $K_{\Z/2}$ 的差别正在于此。 \textit{（实空间的范畴）}\par\smallskip

\noindent \textbf{（原书 p.166）}\ 上面的 $\phi_x\in\Hom_\C(E_{\bar x},F_{\bar x})$ 与紧接的 $\overline{\phi_x}(u)=\overline{\phi_x(\bar u)}$（$u\in E_{\bar x}$）不相容：右式要求 $\phi_x$ 定义在 $E_x$ 上，故第一个 $\Hom_\C$ 的纤维应为 $E_x,F_x$，页图上 $x$ 的上横线疑为误排。 \textit{（实空间的范畴）}\par\smallskip

\noindent \textbf{（原书 p.171）}\ 上面三元组的第一个空间页图作 $X\times B^{p,q}$，而同组的另两个空间与整个完全序列都只用到 $KR^{p,0}$、$KR^{p,1}$，此处应为 $X\times B^{p,0}$，疑为笔误。 \textit{（周期性定理）}\par\smallskip

\noindent \textbf{（原书 p.176）}\ 按 p.175 的图式，虚线对角 $\phi=\chi\theta\beta^p$ 是从 $KR^{p,q}(X)$ 到 $KR^{p,p+q}(X)$ 的复合；正文却写成 $KR^{p,p}(*)\to KR^{0,p}(*)$，把 $X$ 取成一点时下标 $p,q$ 如何对应，须按 C.3 的系数理论自行补足。又，此处的 $b^p$ 与 p.175 的 $c_p=\mu^*(b^{p,1})$ 中的 $b^{p,1}$ 关系原书未明说，此处照抄存疑。 \textit{（系数理论）}\par\smallskip

\noindent \textbf{（原书 p.177）}\ 这一步是本节的关键：$X\times S^{2,0}_+$ 上的复丛 $F_+$ 配上 $\theta$，就等价于给 $X$ 上的丛 $E$ 一个与对合 $\tau$ 相容的同构 $\alpha\colon E\to\tau^*(\overline{E})$。注意 $\alpha$ 的反线性来自复共轭，而 $\theta$ 的同伦类才是唯一被记住的不变量，所以“实丛的同构类”与“自共轭丛的同伦类”才一一对应。 \textit{（系数理论）}\par\smallskip

\noindent \textbf{（原书 p.180）}\ 系 C.3.8 的正合列 $0\to KR^{-q}(X)\to KR^{-q}(X\times S^{p,0})\to KR^{p+1-q}(X)\to 0$ 在 $p\ge3$ 时分裂，本质原因是 $\eta^3=0$（引理 C.3.7）：正合列由乘以 $\eta^p$ 给出，而 $\eta^3=0$ 使得再乘一次就消失。$p=3$ 的情形用 $S^{3,0}\subset P(\C^3)$ 的二次曲面结构及 C.5 节的左逆命题处理。 \textit{（系数理论）}\par\smallskip

\noindent \textbf{（原书 p.183）}\ 新对合 $\tilde m=-w\bar m$ 的要点是 $w^2=1$、$\bar w=w$ 且 $wz=\bar zw$（$w\in\operatorname{Cliff}(\R^{4,4})$ 为体积元）。这三条正好使 $z\tilde m=\widetilde{zm}$，即新对合与 $\C^4=\R^{4,4}$ 的模结构相容；符号取负是为了让 $\bar{\tilde m}=m$（对合性）。$\Lambda^*(\C^4)$ 在旧结构下是 $C_4$ 模、在新结构下是 $C_8$ 模，这正是维数 4 与 8 的跳跃所在。 \textit{（与 Clifford 代数的关系）}\par\smallskip

\section*{监译者解说}
\addcontentsline{toc}{section}{监译者解说}

\noindent \textbf{（原书 p.199）}\ 本节致谢名单的日文名字 OCR 有若干不可靠处（如「五味消紀」「森晃紀」），译文按 K 理论与相关领域通行写法录为「五味清纪」「森晃纪」；若与原书不符，以原书扫描页为准。 \textit{（3. 正文补注）}\par\smallskip
